\chapter{Disaster Recovery and Regulatory Requirements}\label{ch:nw:disaster-recovery-and-regulatory-requirements}

When Superstorm Sandy reached New York in October 2012, the U.S. national securities exchanges ``jointly decided not to open
for trading on October 29 and October 30''. The securities industry's annual test of trading from backup sites had run on
October 27, two days before, and found nothing that would have stopped the markets from opening on their backup systems; many
firms, the SEC later wrote, ``remained uncomfortable about switching over''. Two years later Regulation SCI required the
exchanges' recovery plans to be reasonably designed to resume their critical systems within two hours of a wide-scale
disruption, and required the members each venue designates to test its backup site with it at least once a year.

Book~13 described failover, primary--backup replication and the split brain; Book~6 operational risk; Book~11 the kill
switch; Book~15 incident response. This chapter is about sites and rules: where a firm puts its recovery site, how exchanges
test theirs with their members, which resilience rules a trading firm is examined on, and how a recovery design turns into a
recovery time and a recovery point. The recovery simulation is labelled as such; the rules are quoted.

\section{Secondary sites and how they are placed}

\begin{definition}[Primary data centre, disaster-recovery site]\label{def:nw:disaster-recovery-and-regulatory-requirements:sites}
A firm's \emph{primary data centre}\index{primary data centre} is the site from which it normally trades; its
\emph{disaster-recovery site}\index{disaster-recovery site} is a second site, far enough away not to share the primary's likely
disasters, from which the firm can resume trading when the primary is lost.
\end{definition}

Two distances pull against each other. The recovery site must be outside the hazards that could take the primary, a storm, a
flood, a regional power failure; and it must be close enough for its replication to be cheap and its connections to the venues
usable. The exchanges' own choices show the scale: CME's U.S. BrokerTec markets run in Secaucus and recover in Aurora, its
European ones run in Slough and recover in Secaucus; LSEG's spot matching runs in Slough and recovers in Secaucus
(chapters~23 to~25). Recovery across a continent or an ocean, not across the river.

\begin{table}[ht]
\centering\small
\begin{tabular}{@{}lrl@{}}
\toprule
Candidate recovery site & From Secaucus NY5 (km) & Outside a 100 km hazard radius \\
\midrule
Secaucus NY4 (same campus) & 0.3 & no \\
Carteret & 26.0 & no \\
Mahwah & 34.2 & no \\
Aurora & 1\,191.1 & yes \\
\bottomrule
\end{tabular}
\caption{Candidate recovery sites for a firm trading from Secaucus NY5: geodesic separation (chapter~10's coordinates) against a
hazard radius of \qty{100}{\kilo\metre} (an assumption). Data: \texttt{nw\_dr.separations()}.}\label{tab:nw:disaster-recovery-and-regulatory-requirements:sites}
\end{table}

The New Jersey sites are all inside any storm that closes Secaucus. Aurora is not, and costs what distance costs: synchronous
replication to it would add \qty{17.4}{\milli\second} of round trip to every write at a \omterm{def:nw:the-north-american-map:factor}{route factor} of 1.5, against
\qty{0.5}{\milli\second} to Mahwah, so the firm replicates asynchronously and accepts a recovery point.

\section{Exchange failover tests}

\begin{definition}[Business continuity plan, exchange failover test]\label{def:nw:disaster-recovery-and-regulatory-requirements:test}
A \emph{business continuity plan}\index{business continuity plan} is a firm's documented set of arrangements, people and
procedures for continuing or resuming its critical activities after a disruption. An \emph{exchange failover test}\index{exchange
failover test} is a scheduled exercise in which an exchange runs its markets from its recovery site or backup components and its
members connect, trade and verify from theirs.
\end{definition}

Tests exist because Sandy showed that a tested backup is not a trusted one. Regulation SCI makes them mandatory for the members an
exchange designates, at least once every twelve months, coordinated across the industry. The futures industry runs its own:
FIA's 2025 test, on Saturday 25 October with a pre-test in September, had more than a hundred organisations verify round-trip
connectivity and process recovery between primary and backup sites (\cref{dat:nw:disaster-recovery-and-regulatory-requirements:tests}).
Exchanges also fail over routinely: CME's \omterm{def:nw:futures-connectivity:msgw}{market segment gateways} switch to their backup components and back in weekly windows.

\begin{dated}{2026-09}{Failover tests and exchange recovery sites}\label{dat:nw:disaster-recovery-and-regulatory-requirements:tests}
FIA disaster recovery test: 25 October 2025 (pre-test 27 September 2025), more than 100 participants. CME \omterm{def:nw:futures-connectivity:msgw}{market segment gateways}:
weekly failover windows to backup components. BrokerTec U.S.: NY5, recovering in Aurora, 4-hour \omterm{def:nw:disaster-recovery-and-regulatory-requirements:objectives}{recovery time objective}; BrokerTec
EU: LD4.2, recovering in NY5, 2 hours. LSEG spot matching: LD4, recovering in NY6.
\end{dated}

\section{The resilience rules firms are examined on}

\begin{definition}[Operational resilience, important business service, impact tolerance]\label{def:nw:disaster-recovery-and-regulatory-requirements:resilience}
\emph{Operational resilience}\index{operational resilience} is a firm's ability to prevent, adapt to, respond to, recover from and
learn from operational disruptions. An \emph{important business service}\index{important business service} is a service the firm
provides to clients or the market whose disruption could cause them intolerable harm or threaten the firm or market integrity; its
\emph{impact tolerance}\index{impact tolerance} is the maximum tolerable disruption to it, set by the firm, usually as a time.
\end{definition}

\begin{dated}{2026-09}{Resilience rules, from the regulators' texts}\label{dat:nw:disaster-recovery-and-regulatory-requirements:rules}
\begin{itemize}
\item Regulation SCI (U.S. exchanges and other SCI entities): backup capabilities ``sufficiently resilient and geographically diverse''
  and ``reasonably designed to achieve next business day resumption of trading and two-hour resumption of critical SCI systems
  following a wide-scale disruption''; designated members test with the entity at least once every 12 months.
\item RTS 6, Article 14 (algorithmic trading firms in the EU): documented business continuity arrangements per venue, scenarios
  including the loss of data centres and suppliers, relocation to a back-up site where appropriate, orderly shutdown; reviewed and
  tested annually.
\item FCA PS21/3 (UK): \omterm{def:nw:disaster-recovery-and-regulatory-requirements:resilience}{important business services} and \omterm{def:nw:disaster-recovery-and-regulatory-requirements:resilience}{impact tolerances} set by 31 March 2022; able to remain within them by
  31 March 2025. DORA (EU, Regulation 2022/2554): applies from 17 January 2025.
\end{itemize}
\end{dated}

The rules differ by who they bind but share a shape. The exchange must resume its critical systems within two hours; the trading firm
must have arrangements to resume or shut down in order, tested every year; and in the United Kingdom the firm itself names the
services that matter and the disruption it will tolerate, then shows it can stay within it. A trading firm's recovery design is
therefore judged on two numbers it sets itself, a recovery time and a recovery point, against the tolerance it declared.

\section{Recovery objectives and runbooks}

\begin{definition}[Recovery time objective, recovery point objective]\label{def:nw:disaster-recovery-and-regulatory-requirements:objectives}
A \emph{recovery time objective}\index{recovery time objective} is the longest a system may take, from its loss, to be back in
service; a \emph{recovery point objective}\index{recovery point objective} is the most recent data the firm may lose, measured as the
time between the last change the recovery site holds and the moment of loss.
\end{definition}

A recovery is a runbook: detect the loss and decide to fail over, switch the sessions to the recovery addresses, restore or start what
is not running, load positions and reference data, reconnect and verify sessions, reconcile orders and positions against the
exchanges' drop copies, resume. Each step has a duration that varies; the recovery time is their sum.

\omcode{code/firm/drplan/firm_drplan.py}{61}{84}{Recovery time as the sum of a runbook's steps, the recovery point by replication mode, and the price of synchronous replication.}\label{lst:nw:disaster-recovery-and-regulatory-requirements:rto}

\begin{proposition}[What sets each objective]\label{prop:nw:disaster-recovery-and-regulatory-requirements:objectives}
The recovery point depends only on replication: zero if synchronous, the replication lag plus the one-way delay if asynchronous, up to
the snapshot interval if restored from snapshots. The recovery time depends only on the runbook: it is at least the sum of the steps'
minimum durations, whatever the replication.
\end{proposition}

\begin{proof}
Data the recovery site does not hold when the primary is lost is lost; how much is set by when it was sent, which is the replication.
The runbook's steps run in sequence after the loss, and none can take less than its minimum, so neither objective can be bought with
the other's tool.
\end{proof}

The chapter's firm, trading from Secaucus with a recovery site in Aurora, compares two designs
(\cref{fig:nw:disaster-recovery-and-regulatory-requirements:rto}). The hot design keeps servers running in Aurora with asynchronous
replication (a lag of \qty{50}{\milli\second}) and sessions ready: its runbook is five steps, with medians of 15, 10, 20, 10 and 5 minutes.
The warm design keeps the site and a few servers, restores from snapshots taken every fifteen minutes, and has six steps, including 45
minutes to start and restore and 30 to load data. All durations are assumptions, lognormal around their medians.

\omcode{code/networks/28-disaster-recovery-and-regulatory-requirements/python/nw_dr.py}{24}{34}{The two designs: runbooks, replication and yearly cost.}\label{lst:nw:disaster-recovery-and-regulatory-requirements:plans}

\begin{omfigure}
\begin{tikzpicture}
\begin{axis}[width=11.5cm,height=5.8cm,xmin=0,xmax=300,ymin=0,ymax=0.36,
  xlabel={recovery time (minutes)},ylabel={share of recoveries in each 10-minute bin},
  tick label style={font=\footnotesize},label style={font=\small},grid=major,grid style={gray!25},
  legend pos=north east,legend style={font=\footnotesize},table/col sep=comma]
\addplot[omDef,thick,fill=omDef!25,const plot mark mid,area legend] table[x=minutes,y=hot]{figdata/networks/28-disaster-recovery-and-regulatory-requirements/rto.csv} \closedcycle;
\addplot[omThm,thick,fill=omThm!20,fill opacity=0.6,const plot mark mid,area legend] table[x=minutes,y=warm]{figdata/networks/28-disaster-recovery-and-regulatory-requirements/rto.csv} \closedcycle;
\draw[omMeth,very thick,dashed] (axis cs:120,0) -- (axis cs:120,0.36);
\node[text=omMeth,font=\footnotesize,anchor=west] at (axis cs:122,0.33) {two hours};
\legend{hot site,warm site}
\end{axis}
\end{tikzpicture}

\omcaption{Simulation: the recovery-time distributions of the hot and warm designs, against the two-hour line. The hot design resumes
within two hours in 99.2\% of simulated recoveries (median 64 minutes); the warm one in 20.3\% (median 148). Data:
\texttt{fig\_dr.py}, \texttt{nw\_dr.rto\_hist()}.}\label{fig:nw:disaster-recovery-and-regulatory-requirements:rto}
\end{omfigure}

The hot design resumes within two hours in 99.2\% of simulated recoveries, with a median of 64 minutes and a 95th percentile of 96; its
recovery point is \qty{56}{\milli\second}, the lag plus the one-way delay to Aurora. The warm design meets two hours in 20.3\%, with a
median of 148 minutes and a 95th percentile of 237, and can lose fifteen minutes of data. If the hot site costs 70\% of the primary's
yearly infrastructure and the warm one 25\% (chapter~27's example stack, USD 3.19 million a year; assumptions), the difference is USD
1.44 million a year: the price of the two-hour line for this firm.

\begin{method}[Designing a recovery for a trading firm]\label{met:nw:disaster-recovery-and-regulatory-requirements:design}
\begin{enumerate}
\item Name the \omterm{def:nw:disaster-recovery-and-regulatory-requirements:resilience}{important business services} and set an \omterm{def:nw:disaster-recovery-and-regulatory-requirements:resilience}{impact tolerance} for each; derive the recovery time and point objectives.
\item Choose a recovery site outside the primary's hazards and near the venues' own recovery sites; compute what synchronous
  replication would cost in latency, and replicate asynchronously if it is too much.
\item Write the runbook as timed steps; measure each step in tests, not in estimates.
\item Simulate the recovery time against the objective; buy down the steps that dominate, or change design.
\item Take part in the exchanges' and the industry's failover tests, and run your own at least once a year; record the results as
  evidence for the regulator.
\end{enumerate}
\end{method}

\section{Tutorial: the two-hour rule}\label{tut:nw:disaster-recovery-and-regulatory-requirements:tut}

\begin{tutorial}
\textbf{Goal.} Place a recovery site, simulate the hot and warm designs, and compare them with the two-hour line and their cost.
\textbf{End state:} \cref{fig:nw:disaster-recovery-and-regulatory-requirements:rto} and
\cref{tab:nw:disaster-recovery-and-regulatory-requirements:sites}.
\begin{tutsteps}
\item \textbf{Sites.} \texttt{firm\_drplan.separation\_km} and \texttt{outside\_hazard} on chapter~10's coordinates.
\item \textbf{Plans.} \texttt{Plan} and \texttt{Step} (\cref{lst:nw:disaster-recovery-and-regulatory-requirements:plans});
  \texttt{rto\_samples}, \texttt{p\_within} and \texttt{rpo} (\cref{lst:nw:disaster-recovery-and-regulatory-requirements:rto}).
\item \textbf{Calendar.} \texttt{load\_calendar} and \texttt{tested\_within} for the annual test.
\end{tutsteps}
\textbf{What to change next.} Let some steps run in parallel (reconciliation while sessions reconnect) and see which design's median
moves; add the chance that the decision to fail over is delayed, as it was before Sandy.
\end{tutorial}

\section{Build: the disaster-recovery plan}\label{bld:nw:disaster-recovery-and-regulatory-requirements:drplan}

\begin{build}
\bfield{Purpose} A recovery plan as data with its recovery time and point, checks against targets and a test calendar: the resilience
rows of chapter~29's plan.
\bfield{Interface} \texttt{firm\_drplan}: \texttt{Step}, \texttt{Plan}, \texttt{separation\_km}, \texttt{outside\_hazard},
\texttt{rto\_samples}, \texttt{p\_within}, \texttt{rpo}, \texttt{sync\_penalty\_ms}, \texttt{tested\_within}, \texttt{load\_calendar}.
\bfield{Rules} Sites come from the site table with its sources; test dates are dated rows; runbook durations are stated assumptions,
replaced by measured ones after each test.
\bfield{Acceptance tests} \texttt{code/firm/drplan/tests/}: a deterministic runbook, the recovery point by mode, separations, and the
annual-test check.
\bfield{Stretch} Parallel steps; correlated failures of the primary and a provider; the exchange's own recovery time added to the firm's.
\end{build}

\begin{omsources}
\item 17 CFR 242.1001 and 242.1004 (Regulation SCI); SEC Releases 34-69077 (2013) and 34-73639 (2014).
\item Commission Delegated Regulation (EU) 2017/589, Article 14; FCA PS21/3; ESMA on DORA.
\item FIA, 2025 disaster recovery test.
\end{omsources}

\section{Exercises}

\begin{exercise}[$\star$]\label{exo:nw:disaster-recovery-and-regulatory-requirements:1}
What must an SCI entity's backup capabilities be designed to achieve after a wide-scale disruption?
\end{exercise}

\begin{exercise}[$\star$]\label{exo:nw:disaster-recovery-and-regulatory-requirements:2}
What is the difference between a \omterm{def:nw:disaster-recovery-and-regulatory-requirements:objectives}{recovery time objective} and a \omterm{def:nw:disaster-recovery-and-regulatory-requirements:objectives}{recovery point objective}?
\end{exercise}

\begin{exercise}[$\star$]\label{exo:nw:disaster-recovery-and-regulatory-requirements:3}
Why is Mahwah a poor recovery site for a firm trading from Secaucus?
\end{exercise}

\begin{exercise}[$\star\star$]\label{exo:nw:disaster-recovery-and-regulatory-requirements:4}
Using \cref{prop:nw:disaster-recovery-and-regulatory-requirements:objectives}, why does faster replication not shorten the warm design's
recovery time?
\end{exercise}

\begin{exercise}[$\star\star$]\label{exo:nw:disaster-recovery-and-regulatory-requirements:5}
What does RTS 6 require a trading firm's business continuity arrangements to include about outstanding orders?
\end{exercise}

\begin{exercise}[$\star\star$]\label{exo:nw:disaster-recovery-and-regulatory-requirements:6}
The firm's last failover test was on 25 October 2025. Is it within twelve months on 28 September 2026, and on 1 November 2026?
\end{exercise}

\begin{exercise}[$\star\star\star$]\label{exo:nw:disaster-recovery-and-regulatory-requirements:7}
\emph{Coding.} With \texttt{firm\_drplan}, how much would halving the warm design's restore step (45 to 22.5 minutes) raise its chance of
resuming within two hours?
\end{exercise}

\begin{exercise}[$\star\star\star$]\label{exo:nw:disaster-recovery-and-regulatory-requirements:8}
\emph{Find the flaw.} ``We tested our backup site last month and it worked, so we will switch to it if a storm comes.''
\end{exercise}

\section{Problem: The Two-Hour Rule}

\begin{problem}[{Weekend problem --- hot against warm recovery}]\label{pb:nw:disaster-recovery-and-regulatory-requirements:1}
A firm trades from Secaucus NY5 and must choose a recovery design. It has set itself the exchanges' two-hour line as its \omterm{def:nw:disaster-recovery-and-regulatory-requirements:objectives}{recovery time
objective}. Use the chapter's model.

\medskip\noindent\textbf{Part I --- The site.}
\begin{enumerate}
\item How far is each candidate site from NY5, and which lie outside a 100-kilometre hazard radius?
\item What would synchronous replication to Aurora and to Mahwah add to each write?
\item Why do the exchanges recover so far away?
\item What recovery point does asynchronous replication to Aurora give?
\end{enumerate}

\medskip\noindent\textbf{Part II --- The designs.}
\begin{enumerate}[resume]
\item What is the hot design's chance of resuming within two hours, its median and 95th percentile?
\item And the warm design's?
\item What can the warm design lose in data?
\item Which steps dominate the warm design's time?
\end{enumerate}

\medskip\noindent\textbf{Part III --- The rules.}
\begin{enumerate}[resume]
\item What does Regulation SCI require of exchanges, and of their designated members?
\item What does RTS 6 require of trading firms?
\item What did the UK rules require by 31 March 2025?
\item What did Sandy show that tests alone did not?
\end{enumerate}

\medskip\noindent\textbf{Part IV --- The verdict.}
\begin{enumerate}[resume]
\item State the \emph{named result}: the probability that warm and hot designs resume within two hours, and the annual cost of the
  difference.
\item Which inputs are published and which assumed?
\item Is the hot design worth USD 1.44 million a year?
\item What would make the warm design meet the line more often?
\item What evidence would a regulator ask for?
\item How should the firm take part in exchange tests?
\item Where does this go in chapter~29's plan?
\item In one sentence: what does a recovery design buy?
\end{enumerate}
\end{problem}

\section{Interview questions}

\begin{interviewq}[$\star$ \iqroles{developer}]\label{iq:nw:disaster-recovery-and-regulatory-requirements:1}
What are RTO and RPO, and what sets each?
\end{interviewq}

\begin{interviewq}[$\star\star$ \iqroles{developer}]\label{iq:nw:disaster-recovery-and-regulatory-requirements:2}
Where would you put the recovery site of a firm trading in Secaucus, and why?
\end{interviewq}

\begin{interviewq}[$\star\star$ \iqroles{developer}]\label{iq:nw:disaster-recovery-and-regulatory-requirements:3}
How do you avoid a split brain when failing over a trading system?
\end{interviewq}

\begin{interviewq}[$\star\star$ \iqroles{trader}]\label{iq:nw:disaster-recovery-and-regulatory-requirements:4}
The primary site is lost with open positions. What happens in the first thirty minutes?
\end{interviewq}

\begin{interviewq}[$\star\star$ \iqroles{developer, researcher}]\label{iq:nw:disaster-recovery-and-regulatory-requirements:5}
How would you decide between a hot and a warm recovery site?
\end{interviewq}

\begin{interviewq}[$\star\star\star$ \iqroles{developer, researcher}]\label{iq:nw:disaster-recovery-and-regulatory-requirements:6}
Design the business continuity arrangements of an algorithmic trading firm active in the U.S. and the EU, so that they satisfy the rules
it is examined on.
\end{interviewq}
